% Lösungen Einheit 2
\einheitenkopf*{Einheit 2 von 5 · Wachstumsfaktor und Wachstumstabelle}

\erg{21}{a) nimmt zu \quad b) nimmt ab \quad c) nimmt zu \quad d) nimmt ab \quad e) nimmt zu}
\erg{22}{a) $q = 1{,}1$ \quad b) $q = 1{,}15$ \quad c) $q = 1{,}02$ \quad d) $q = 1{,}35$ \quad e) $q = 0{,}6$ \quad f) Abnahme um $7\,\%$ \quad g) $290 : 250 = 1{,}16$ \quad h) $460 : 400 = 529 : 460 = 608{,}35 : 529 = 1{,}15$: Der Wert wächst jedes Jahr mit dem Faktor $1{,}15$, also um $15\,\%$. \quad i) $1\,936 : 1\,600 = 1{,}21 = 1{,}1 \cdot 1{,}1$, also $q = 1{,}1$ und $10\,\%$ \quad j) $320 : 256 = 1{,}25$; $400 : 320 = 1{,}25$; $500 : 400 = 1{,}25$ – alle Quotienten gleich, Faktor $1{,}25$}
\erg{23}{a) $140$ \quad b) $550$ \quad c) $840$ \quad d) $255$ \quad e) $400;\ 440;\ 484$ (der Anfangswert steht im Jahr 0) \quad f) $1\,650$ \quad g) $3\,000 \cdot 1{,}05^3 \approx 3\,472{,}88\,€$}
\erg{24}{a) $1\,331 : 1{,}1 = 1\,210$ Kaninchen \quad b) Jahr 4: $500 \cdot 1{,}1^4 = 732{,}05$ \quad c) $480;\ 384$ \quad d) 2025: $144{,}00\,€$; 2027: $149{,}04 \cdot 1{,}035 \approx 154{,}26\,€$}
\erg{25}{a) Fehler: Paul zieht jedes Jahr denselben Betrag ($1\,080\,€$) ab; richtig ist jedes Jahr mal $0{,}85$: $7\,200;\ 6\,120;\ 5\,202;\ 4\,421{,}70\,€$ \quad b) $3\,240;\ 2\,916;\ 2\,624{,}40\,€$ \quad c) Fehler: Lena hat nur mit $0{,}09$ malgenommen und damit nur die Erhöhung berechnet; der Faktor ist $1{,}09$: $2\,000 \cdot 1{,}09 = 2\,180\,€$ \quad d) $2\,500 \cdot 1{,}07 = 2\,675\,€$}
\erg{26}{a) $0{,}9 = 1 - 0{,}1$: vom Wert bleiben $90\,\%$, $10\,\%$ fallen weg. \quad b) Der Quotient sagt, mit welcher Zahl man malnimmt; die Differenz ist nur der Zuwachs in einem Jahr, und der wird bei exponentiellem Wachstum jedes Jahr größer. \quad c) Nein: Das Sinken bezieht sich auf den höheren Preis; $1{,}5 \cdot 0{,}5 = 0{,}75$, der Preis ist danach um $25\,\%$ niedriger.}
\erg{27}{a) Jahr 1: $12\,240$ \quad Jahr 2: $12\,485$ \quad Jahr 3: $12\,734$ \quad b) Ja, ab Jahr 3 ($12\,734 > 12\,700$).}
